\section{Kesimpulan}
\label{sec:conclusion}

SkripsiBench mengukur beberapa hal yang dibutuhkan karya skripsi Indonesia dari sebuah model bahasa besar---kerangka yang sesuai dengan program studi, penggunaan sumber yang setia, rujukan yang nyata, bahasa Indonesia baku, kode analisis yang benar dan LaTeX yang dapat dikompilasi---bersama dengan biaya tiap jawaban dan model yang menghasilkannya. Di antara \NModels{} model, tiga model API terbaik tidak berbeda signifikan sementara biaya model API mencakup dua orde besaran, rujukan paling membedakan model-model tersebut, dan rata-rata model lokal tertinggal pada setiap tugas, paling kecil pada bahasa Indonesia baku. Saran praktisnya adalah memilih berdasarkan harga di antara model-model teratas, memverifikasi setiap rujukan, dan menggunakan model lokal ketika privasi mengharuskannya. Paket ini mengevaluasi model lain dengan satu perintah; tier premium, reasoning effort tinggi, pelaksanaan berulang dan pemeriksaan manusia terhadap para juri dan kunci jawaban adalah langkah selanjutnya.
